\newpage

%========================================================================================
%	附录
%========================================================================================
\begin{appendices}
\section{调查问卷}

\begin{center}
\textbf{杭甬温环境噪声污染调查问卷}
\end{center}

尊敬的受访者：

您好！为精准掌握杭甬温环境噪声污染现状，了解噪声对居民生活的影响，为噪声治理提供科学数据支撑，我们特开展本次问卷调查。调研结果仅用于学术研究与政策参考，问卷填写耗时较短，请您根据实际情况如实作答。衷心感谢您在百忙之中抽出宝贵时间支持本次调查工作！

\begin{center}
\textbf{第一部分　基本信息}
\end{center}

\textbf{1. 您的性别}\\
$\circ$男 \qquad $\circ$女

\textbf{2. 您的年龄}\\
$\circ$18岁及以下 \quad $\circ$18–30岁 \quad $\circ$31–45岁 \quad $\circ$46–60岁 \quad $\circ$61岁及以上

\textbf{3. 您现居住地}\\
$\circ$杭州 \quad $\circ$宁波 \quad $\circ$温州

\textbf{4. 您的居住区域类型}\\
$\circ$城市主城区 \quad $\circ$城市郊区 \quad $\circ$县城 \quad $\circ$乡镇 \quad $\circ$农村

\textbf{5. 您的职业类型}\\
$\circ$机关事业单位工作人员 \quad $\circ$农林牧渔从业者 \quad $\circ$企业工作人员\\
$\circ$自由职业者/个体工商户 \quad $\circ$学生 \quad $\circ$军人/武警\\
$\circ$退休人员 \quad $\circ$无业 \quad $\circ$其他 \underline{\hspace{3cm}}

\textbf{6. 您的文化程度}\\
$\circ$小学及以下 \quad $\circ$初中 \quad $\circ$高中/中专/技校 \quad $\circ$大专 \quad $\circ$本科 \quad $\circ$硕士及以上

\textbf{7. 您居住房屋墙体是否为隔音墙}\\
$\circ$是 \quad $\circ$否 \quad $\circ$不清楚

\textbf{8. 您是否需要居家办公（仅题5未选学生、退休、无业的对象回答）}\\
$\circ$长期居家办公 \quad $\circ$偶尔居家办公 \quad $\circ$从不居家办公

\begin{center}
\textbf{第二部分　噪声暴露与环境感知}
\end{center}

\textbf{9. 您日常生活中受到环境噪声的影响大吗？}\\
$\circ$严重影响生活 \quad $\circ$影响较大 \quad $\circ$影响一般 \quad $\circ$偶尔有影响 \quad $\circ$几乎没有影响

\textbf{10. 您的居住区域隔音水平如何？}\\
$\circ$很好 \quad $\circ$较好 \quad $\circ$一般 \quad $\circ$较差 \quad $\circ$很差

\textbf{11. 各类噪声对您的干扰程度（请在对应选项处打勾）}

\begin{table}[H]
\centering
\rowcolors{2}{white}{white}
\begin{tabular}{|M{5.2cm}|M{1.6cm}|M{1.6cm}|M{1.6cm}|M{1.6cm}|M{1.6cm}|}
\hline
\rowcolor{tabhead}
\thead{噪声类型} & \thead{无干扰} & \thead{轻微干扰} & \thead{一般干扰} & \thead{较大干扰} & \thead{严重干扰} \\
\hline
交通噪声（车流、鸣笛等） & & & & & \\
\hline
周围人生活噪声（走动、打游戏、交谈等） & & & & & \\
\hline
装修/施工噪声（地铁施工等） & & & & & \\
\hline
商业经营噪声（KTV、摆摊等） & & & & & \\
\hline
公共活动噪声（广场舞、聚众娱乐等） & & & & & \\
\hline
动物噪声（犬吠、猫叫等） & & & & & \\
\hline
\end{tabular}
\end{table}

\textbf{12. 您离噪声源近吗？}\\
$\circ$周围无明显噪声源 \quad $\circ$距离较远（$<$100m） \quad $\circ$距离适中（$<$50m）\\
$\circ$距离较近（$<$20m） \quad $\circ$近在咫尺（$<$10m）

\textbf{13. 哪个时段的噪音对你影响最大？}\\
$\circ$上午（6:00-12:00） \quad $\circ$下午（12:00-18:00） \quad $\circ$晚上（18:00-22:00） \quad $\circ$深夜（22:00-次日6:00）

\begin{center}
\textbf{第三部分　噪声负面影响（请在对应选项处打勾）}
\end{center}

\begin{table}[H]
\centering
\rowcolors{2}{white}{white}
\begin{tabular}{|M{5.2cm}|M{1.6cm}|M{1.6cm}|M{1.6cm}|M{1.6cm}|M{1.6cm}|}
\hline
\rowcolor{tabhead}
\thead{调研问题} & \thead{非常不认同} & \thead{比较不认同} & \thead{一般} & \thead{比较认同} & \thead{非常认同} \\
\hline
噪声会延长我的入睡时间，使我难以快速入眠 & & & & & \\
\hline
噪声会让我半夜惊醒、中断睡眠 & & & & & \\
\hline
噪声会影响我的注意力，导致无法专心做事 & & & & & \\
\hline
噪声会损害我的邻里、室友、家庭等关系 & & & & & \\
\hline
噪声会让我变得更加疲惫、精神状态变差 & & & & & \\
\hline
噪声导致我产生听觉不适（如耳鸣、听觉疲劳等） & & & & & \\
\hline
噪声导致我产生非听觉类的身体不适（心悸、月经失调等等） & & & & & \\
\hline
噪声会让我变得冲动、易怒、烦躁不安 & & & & & \\
\hline
\end{tabular}
\end{table}

\begin{center}
\textbf{第四部分　治理评价与应对方法（请在对应选项处打勾）}
\end{center}

\begin{table}[H]
\centering
\rowcolors{2}{white}{white}
\begin{tabular}{|M{5.2cm}|M{1.6cm}|M{1.6cm}|M{1.6cm}|M{1.6cm}|M{1.6cm}|}
\hline
\rowcolor{tabhead}
\thead{调研问题} & \thead{非常不满意} & \thead{不太满意} & \thead{一般} & \thead{比较满意} & \thead{非常满意} \\
\hline
您对居住地周边整体声环境质量的满意度 & & & & & \\
\hline
您对本地现阶段噪声治理工作整体满意度 & & & & & \\
\hline
\end{tabular}
\end{table}

\textbf{24. 当遭遇噪声污染扰民时，您的主要应对方式：}\\
$\square$默默忍受，不采取任何措施\\
$\square$自行与噪声产生方沟通协商\\
$\square$向物业、社区居委会反馈投诉\\
$\square$佩戴耳塞、加装隔音门窗等自我防护\\
$\square$拨打政务热线（12345）、公安、环保等官方渠道投诉\\
$\square$采取法律途径维权\\
$\square$其他 \underline{\hspace{3cm}}

\begin{center}
\textbf{第五部分　噪声治理认知（请在对应选项处打勾）}
\end{center}

\begin{table}[H]
\centering
\rowcolors{2}{white}{white}
\begin{tabular}{|M{5.2cm}|M{1.6cm}|M{1.6cm}|M{1.6cm}|M{1.6cm}|M{1.6cm}|}
\hline
\rowcolor{tabhead}
\thead{调研问题} & \thead{非常不了解} & \thead{不太了解} & \thead{一般} & \thead{比较了解} & \thead{非常了解} \\
\hline
您对噪声污染投诉、维权渠道的了解程度（12345政务热线、浙里办等） & & & & & \\
\hline
您对《中华人民共和国噪声污染防治法》的了解程度 & & & & & \\
\hline
\end{tabular}
\end{table}

\textbf{27. 您认为是否有必要强化浙江省环境噪声治理？}\\
$\circ$非常有必要 \quad $\circ$有必要 \quad $\circ$一般 \quad $\circ$没有必要 \quad $\circ$完全没有必要

\textbf{28. 您认为环境噪声治理的责任主体有哪些？}\\
$\square$政府生态环境、城管、公安等职能部门\\
$\square$小区物业、社区居委会\\
$\square$施工单位、商铺企业等噪声产生主体\\
$\square$全体居民（自觉遵守、互相约束）\\
$\square$其他 \underline{\hspace{3cm}}

\begin{center}
\textbf{第六部分　噪声治理建议}
\end{center}

\textbf{29. 您认为改善噪声污染需要哪些举措：}\\
$\square$加大监管巡查与执法处罚力度\\
$\square$完善管控法律法规与区域管控标准\\
$\square$增设隔音屏、隔音墙体等硬件设施\\
$\square$合理规范施工、交通等的噪声排放时段\\
$\square$加强文明宣传，增强居民自律意识\\
$\square$完善噪声投诉处理机制，提高办事效率\\
$\square$优化城市声环境规划（如高架桥远离学校等）\\
$\square$其他 \underline{\hspace{3cm}}

\textbf{30. 对于浙江省环境噪声污染治理，您有哪些具体意见与建议？}\\
\underline{\hspace{\textwidth}}\\
\underline{\hspace{\textwidth}}\\
\underline{\hspace{\textwidth}}

\section{访谈记录表}

\subsection{访谈人物1}

\textbf{一、访谈基本信息}

访谈对象：张阿姨（化名），68岁，独居，现居住于宁波市鄞州区老旧小区

访谈时间：2026年5月1日 上午9:30-10:00

访谈人员：调研人员（以下简称“调研”）

访谈目的：了解主城区独居老人的环境噪声暴露情况、噪声影响及对噪声治理的认知与建议，补充问卷调查数据，精准掌握杭甬温主城区噪声污染现状

\textbf{二、访谈内容实录}

\textbf{调研}：张阿姨，您好！非常感谢您愿意抽出时间参与我们的环境噪声污染调研，我们今天主要想了解一下您日常生活中遇到的噪声情况，您根据自己的实际感受回答就好，不用有顾虑。

\textbf{张阿姨}：没事没事，我一个人在家也清闲，能帮上忙就好，你们尽管问。

\textbf{调研}：首先想了解下您的基本情况，您在这个小区住多久了？房屋墙体是隔音墙吗？

\textbf{张阿姨}：住了快20年了，老小区了，墙体就是普通的砖墙，应该不是隔音墙吧，我也不太懂这些，反正隔壁说话声音大一点能听见，但也不算特别吵。我一个人住，孩子都在外地，平时就自己在家看看电视、喝喝茶。

\textbf{调研}：那您日常生活中，受到环境噪声的影响大吗？有没有哪种噪声会让您觉得不舒服？

\textbf{张阿姨}：影响不算大，就是一般吧。我们这平时会有汽车声，喇叭声，在上下班高峰期，声音能明显听到，但也不至于影响正常生活。还有就是偶尔楼上小孩蹦蹦跳跳，声音有点大，但不会持续很久，忍一忍也就过去了。

\textbf{调研}：那像周围邻居的生活噪声，比如走动、交谈，还有小区里的广场舞、犬吠这些，对您有干扰吗？

\textbf{张阿姨}：邻居都挺客气的，走动、交谈的声音不大，我平时看电视开大点声音，基本听不到。广场舞就小广场，晚上7点多到8点多，离我们这有点距离，不影响我休息。

\textbf{调研}：您觉得居住的区域隔音水平怎么样？平时离这些噪声源近吗？哪个时段的噪声对您影响相对大一点？

\textbf{张阿姨}：隔音就一般吧，老房子了，不能要求太高。楼上要是吵架，能偶尔听见几句，但不频繁。噪声源嘛，马路离小区也就几十米，不算太远也不算太近。时段的话，下班高峰期车流声大一点。还有就是偶尔周末上午，有人装修。其他时段都还好，深夜基本没什么噪声，不影响我睡觉。

\textbf{调研}：那噪声有没有影响到您的睡眠、精神状态这些呢？比如难以入睡、半夜惊醒，或者觉得疲惫、烦躁？

\textbf{张阿姨}：没有没有，我睡眠还挺好的，晚上一关窗，基本听不到什么声音，不会半夜被吵醒，没有不舒服。

\textbf{调研}：那您对现在居住地周边的声环境质量，还有本地的噪声治理工作，满意吗？

\textbf{张阿姨}：挺满意的，毕竟是老小区，人多车多，能做到这样已经不错了。平时也没怎么看到特别吵的情况，偶尔有装修太吵的，物业也会去说，让他们注意时间。噪声治理这块，我觉得政府也挺用心的。

\textbf{调研}：您了解噪声污染的投诉渠道，比如12345政务热线、浙里办这些吗？还有《中华人民共和国噪声污染防治法》，您听说过吗？

\textbf{张阿姨}：12345知道，平时看电视、听广播会提到，但没打过，也不知道具体怎么用。那个什么噪声防治法，没听说过，不太了解这些法律规定。年纪大了，我也不太关注这些。

\textbf{调研}：那您觉得有必要强化浙江省的环境噪声治理吗？您觉得噪声治理应该是谁的责任呢？

\textbf{张阿姨}：还是很有必要的，虽然现在噪声不算大，但能治理得更好一点就更好了。责任的话，我觉得政府肯定要负责，还有那些施工的、开店的，他们不能随便制造噪声。我们居民也得注意，大家互相体谅就好。

\textbf{调研}：那您觉得改善噪声污染，需要做些什么呢？还有没有什么具体的意见或建议？

\textbf{张阿姨}：我觉得吧，首先马路上可以多装几个隔音屏，尤其是靠近小区的路段，减少车流声。还有就是多宣传宣传，让大家都注意点，邻里之间互相理解。其他的也没什么特别的建议，现在这样就挺好，只要能保持住，再稍微改善一下就更好了。

\textbf{调研}：非常感谢张阿姨的配合，您说的这些对我们的调研非常有帮助，再次感谢您支持我们的工作！

\subsection{访谈人物2}

\textbf{一、访谈基本信息}

访谈对象：李先生（化名），45岁，已婚，现居住于温州市乐清市郊区工厂周边职工宿舍，系附近机械零件制造加工厂一线操作工

访谈时间：2026年5月5日 下午19:00-19:30

访谈人员：调研人员（以下简称“调研”）

访谈目的：了解郊区工厂旁职工的环境噪声暴露情况、噪声严重影响及对噪声治理的认知与建议，补充问卷调查数据，精准掌握杭甬温郊区工厂周边噪声污染现状

\textbf{二、访谈内容实录}

\textbf{调研}：李先生，您好！感谢您在工作间隙抽出时间参与我们的环境噪声污染调研，我们主要想了解您居住在工厂周边，日常受到的噪声影响情况，麻烦您根据实际感受如实回答就好。

\textbf{李先生}：没事，反正下班了也没什么事，你们问吧。

\textbf{调研}：首先想了解下您的基本情况，您在这个职工宿舍住多久了？宿舍墙体是隔音墙吗？平时需要居家办公吗？

\textbf{李先生}：住了快5年了，这里是工厂统一的职工宿舍，墙体就是普通的水泥墙，肯定不是隔音墙，特别薄。我是一线操作工，不用居家办公，平时上班8小时，下班就回宿舍休息，偶尔休息的时候在家待着，噪声特别影响人。我们宿舍离工厂也就不到50米，工厂的机器声一天到晚不停，特别吵。

\textbf{调研}：那您日常生活中，受到环境噪声的影响大吗？主要是哪种噪声对您干扰最严重？

\textbf{李先生}：影响特别大，主要就是旁边机械加工厂的机器噪声，那种金属切割、机器运转的轰鸣声，特别刺耳，白天晚上都有。工厂有时候赶工期，深夜也不停工，噪声一点都不减弱，有时候还有货车进出工厂的鸣笛声、装卸货物的碰撞声，确实没法好好休息。除了工厂噪声，偶尔有周边村民的生活噪声，但跟工厂噪声比起来，根本不算什么，主要还是工厂的噪声最扰民。

\textbf{调研}：那像装修/施工噪声、商业经营噪声、公共活动噪声这些，对您有干扰吗？您觉得居住区域的隔音水平怎么样？

\textbf{李先生}：装修噪声很少，周边都是工厂和职工宿舍，基本没人装修；商业经营噪声也少，郊区没什么商铺、KTV，就几个小超市，声音不大。公共活动噪声也几乎没有，这边都是工厂职工，大家下班都比较累，很少有聚众娱乐的情况。隔音水平特别差，我刚才也说了，宿舍墙体薄，不是隔音墙，工厂的机器声隔着墙都能听得清清楚楚，有时候关紧窗户、拉上窗帘，也没什么用，噪声还是能钻进来，隔壁宿舍说话声音大一点都能听见，更别说工厂的轰鸣声了。

\textbf{调研}：您离噪声源（工厂）这么近，哪个时段的噪声对您影响最大？噪声有没有影响到您的睡眠、精神状态这些？

\textbf{李先生}：深夜的时候影响最大，本来上班一天就累，想好好睡一觉，结果工厂深夜还在开工，机器轰鸣声一直响，特别刺耳，我经常难以入睡，有时候好不容易睡着了，还会被突然变大的噪声惊醒，一夜醒好几次，第二天上班都没精神，特别疲惫。有时候白天休息，想在家安静待一会儿、看看手机，噪声也一直吵着，根本静不下来，注意力都集中不了。

\textbf{调研}：那噪声有没有让您出现听觉不适或者其他身体不适的情况？

\textbf{李先生}：有一点，长期听这种刺耳的噪声，我现在偶尔会耳鸣，有时候下班回到宿舍，耳边还会有机器运转的嗡嗡声，得缓好长时间才能消失。有时候还会觉得心悸、头晕，尤其是熬夜被噪声吵醒之后，胸口特别闷。

\textbf{调研}：当遭遇这种严重的噪声污染扰民时，您和身边的同事主要会采取什么应对方式？

\textbf{李先生}：一开始我们都是默默忍受，觉得忍忍就过去了，后来噪声越来越严重，影响到正常休息和身体健康，我们就一起找工厂负责人协商，但是没什么用。也有人戴耳塞、加装隔音窗帘，但是效果也都不好，那种刺耳的机器声还是能听到。也想过打12345政务热线投诉，但是不知道具体怎么说，身边也没人打过，还有人说投诉也没用，我们也没什么更好的办法，只能一直忍，实在受不了了，就临时去外面找个宾馆住一晚。

\textbf{调研}：您了解噪声污染的投诉渠道，比如12345政务热线、浙里办这些吗？还有《中华人民共和国噪声污染防治法》，您听说过吗？

\textbf{李先生}：听说过12345政务热线，但是不知道具体怎么投诉，也不知道投诉之后能不能解决，所以一直没打。不知道我们这种情况，能不能通过法律途径维权，毕竟我们都是普通职工，也没什么文化，不懂这些法律知识。

\textbf{调研}：那您觉得有必要强化浙江省的环境噪声治理吗？您觉得噪声治理的责任主体有哪些？

\textbf{李先生}：太有必要了！再不管的话，我们的身体健康都要出大问题了。责任主体肯定是工厂啊，他们是噪声的产生者，就应该负责控制噪声，应该安装隔音设施，规范生产时间，深夜不能开工。然后政府相关部门，比如什么环保、城管，也应该负责监管，定期去工厂检查，对违规制造噪声的工厂进行处罚，不能放任不管。还有社区居委会，也应该帮我们反映问题，协调我们和工厂之间的矛盾，而不是敷衍我们。

\textbf{调研}：那您觉得改善噪声污染，需要做些什么呢？还有没有什么具体的意见或建议？

\textbf{李先生}：我觉得首先，政府部门要加大监管和执法力度，定期检查工厂的噪声排放情况，对超标排放噪声、深夜违规生产的工厂，要严厉处罚，责令他们整改，控制噪声排放。然后，要让我们投诉之后，能有人及时处理。另外，也可以给我们职工宿舍加装隔音设施，比如隔音墙体、隔音窗户，减少噪声的影响。

\textbf{调研}：非常感谢李先生的坦诚分享，您说的这些对我们的调研非常有价值，也能帮助我们更好地了解郊区工厂周边的噪声污染现状，再次感谢您的配合！

\section{调研过程记录}

\begin{figure}[H]
\centering
\begin{subfigure}[b]{0.32\textwidth}
  \centering
  \includegraphics[width=\textwidth]{图片22.jpg}
  \caption{小组调研前期准备}
\end{subfigure}
\hfill
\begin{subfigure}[b]{0.32\textwidth}
  \centering
  \includegraphics[width=\textwidth]{图片23.jpg}
  \caption{宁波市鄞州区某老旧小区}
\end{subfigure}
\hfill
\begin{subfigure}[b]{0.32\textwidth}
  \centering
  \includegraphics[width=\textwidth]{图片24.jpg}
  \caption{老旧小区独居老人}
\end{subfigure}

\vspace{0.8em}

\hspace*{\fill}
\begin{subfigure}[b]{0.32\textwidth}
  \centering
  \includegraphics[width=\textwidth]{图片25.jpg}
  \caption{温州市某工厂职工}
\end{subfigure}
\hspace{0.02\textwidth}
\begin{subfigure}[b]{0.32\textwidth}
  \centering
  \includegraphics[width=\textwidth]{图片26.jpg}
  \caption{宁波市某校边店老板}
\end{subfigure}
\hspace*{\fill}
\caption{调研过程记录}
\end{figure}

\section{核心代码}

本文的聚类分析与机器学习模型对比采用 Python 实现；游程检验、Spearman相关分析、信度与效度检验以及有序Logistic回归的参数估计依托 SPSS 完成，其操作过程与输出结果已在正文相应章节说明，此处不再重复列出代码。以 Python 实现的核心代码文件清单如下：

\begin{table}[H]
\centering
\caption{核心代码文件清单}
\rowcolors{2}{tablight}{white}
\begin{tabular}{M{6cm}M{8cm}}
\bluetoprule
\rowcolor{tabhead}
\thead{文件名} & \thead{功能说明} \\
\hline
preprocess.py & 数据加载、清洗、编码与特征构造 \\
q3\_cluster.py & KMeans++聚类与人群画像 \\
q5\_compare.py & 有序Logistic与随机森林、XGBoost的交叉验证对比 \\
q6\_tune.py & XGBoost贝叶斯调参 \\
\bluebottomrule
\end{tabular}
\end{table}

以下给出以 Python 实现的关键代码。

\subsection{聚类分析}

\lstinputlisting[language=python]{code/model/q3_cluster.py}

\subsection{机器学习模型对比与调参}

\lstinputlisting[language=python]{code/model/q5_compare.py}

\lstinputlisting[language=python]{code/model/q6_tune.py}

\end{appendices}
